% =================================================================
\section{Introduction}
\label{sec:intro}
% =================================================================

Mosquitoes transmit diseases that kill hundreds of thousands of people
every year.
Knowing which species are present in a region---and in what numbers---is
the foundation of any targeted control effort, but traditional
identification requires physical capture and microscopy.
Miniaturised microphones can passively record the characteristic
wingbeat sound of flying mosquitoes, and machine learning can identify
the species from that sound alone~\cite{mukundarajan2017,kiskin2020humming}.
The practical barrier is \emph{domain shift}: a classifier trained on
recordings from one microphone or environment performs poorly when
tested on a different one, even for the same species.

The BioDCASE 2026 Cross-Domain Mosquito Species Classification
challenge~\cite{biodcase2026} formalises this problem.
Nine species are recorded across five conditions (D1--D5), and the
key metric is \emph{species balanced accuracy on unseen conditions}
($\BAu$)---how well the model recognises each species in a recording
environment it has never trained on.
A companion metric, the \emph{domain generalisation gap}
($\DSG{=}|\BAs - \BAu|$, lower is better), quantifies how much worse
the model gets on new conditions compared to known ones.

We diagnose a critical problem with the dataset that has not been
widely acknowledged.
One recording condition---D5---accounts for 212,339 of the 213,647
training clips: \textbf{99.4\%}.
The other four conditions together provide only 1,308 clips.
This extreme imbalance means that during training, a model almost
never sees examples from D1--D4.
Standard domain-generalisation methods like DANN~\cite{dann2016} and
MixStyle~\cite{mixstyle2021} are built on the assumption that training
data covers all domains; here, they do not.
Applied naively, they improve $\BAu$ by only 0.4 and 0.2 percentage
points respectively---essentially nothing.
The official challenge metric is similarly affected: roughly 85\% of
the test set looks like D5 by recording characteristics, so the
reported $\BAu{=}0.175$ mostly reflects D5 generalisation, not
genuine cross-condition transfer.
We adopt a Leave-One-Domain-Out~(LODO) evaluation instead, holding
out each of D1--D4 in turn, to measure true cross-condition
performance (Section~\ref{sec:data}).

Our contributions:
\begin{itemize}
  \item \textbf{A diagnostic finding:} We show that fixing the domain
    imbalance is a \emph{prerequisite} for domain-generalisation
    methods to work at all.
    Domain-balanced resampling alone gains +15.6~pp $\BAu$---the
    single largest improvement in our ablation.
  \item \textbf{FBS-Mix:} A data-motivated augmentation that mixes
    only the recording-noise frequency band
    (mel bins 0--8, $\approx$36--325~Hz) and leaves the wingbeat band
    (bins 9--35, $\approx$361--1{,}360~Hz) untouched, guided by a
    per-frequency variance analysis.
  \item \textbf{A comprehensive ablation:} We test ten method
    combinations under LODO, including actionable negative results:
    test-time batch-norm adaptation (TTBN) is harmful here, unbalanced
    DANN is near-useless, and a transformer backbone (AST) fails to
    benefit from adversarial training because its gradient reversal
    layer never suppresses domain information.
\end{itemize}

Our best system---balanced resampling, domain-adversarial
training~(DANN), and domain-invariant contrastive
learning~(DiCL)~\cite{drbiol2025} with a projection head---achieves
$\BAu{=}0.378$ (+21.3~pp over the unbalanced MTRCNN~\cite{mtrcnn2025}
baseline) and $\DSG{=}0.202$.
FBS-Mix gives a lighter-weight alternative: $\DSG{=}0.229$ with no
model changes.
